% Hart's Theorem
\documentclass{standalone}
\usepackage{ctex}
\usepackage{tikz}

\usetikzlibrary{cartes}

% Styles
\tikzstyle{axes}=[]
\tikzstyle{important line}=[very thick]
\tikzstyle{information text}=[rounded corners,fill=red!10,inner sep=1ex]

\begin{document}

\begin{tikzpicture}
  \tikzmath{
    \A{1} = 0; \A{2} = 3; 
    \B{1} = -1; \B{2} = -2;
    \C{1} = 3; \C{2} = -2;
    \kd = -0.9;
    \ke = 0.25;
  }

  \tikzset{
    plane/combine={\B,\C,\kd,1-\kd,\D},
    plane/combine={\A,\C,\ke,1-\ke,\E},
    plane/intersect ll={\A,\B,\D,\E,\F},
    plane/circumcenter={\A,\B,\C,\Oa},
    plane/circumradius={\A,\B,\C,\ra},
    plane/circumcenter={\A,\E,\F,\Ob},
    plane/circumradius={\A,\E,\F,\rb},
    plane/circumcenter={\D,\B,\F,\Oc},
    plane/circumcenter={\D,\C,\E,\Od},
    plane/intersect cc={\Oa,\ra,\Ob,\rb,\P,\Q},
    plane/exclude={\P,\Q,\A,\M},
    plane/incenter={\A,\B,\C,\Ia},
    plane/incenter={\A,\E,\F,\Ib},
    plane/incenter={\D,\B,\F,\Ic},
    plane/incenter={\D,\C,\E,\Id},
    plane/excenter={\A,\B,\C,\Ja},
    plane/excenter={\A,\E,\F,\Jb},
    plane/excenter={\D,\B,\F,\Jc},
    plane/excenter={\D,\C,\E,\Jd},
    plane/excenter={\B,\C,\A,\Ka},
    plane/excenter={\E,\F,\A,\Kb},
    plane/excenter={\B,\F,\D,\Kc},
    plane/excenter={\C,\E,\D,\Kd},
    plane/excenter={\C,\A,\B,\La},
    plane/excenter={\F,\A,\E,\Lb},
    plane/excenter={\F,\D,\B,\Lc},
    plane/excenter={\E,\D,\C,\Ld},
  }

  % \axes[grid] (-4:8,-4:4);
  \coordinate (A) at (\A{1},\A{2});
  \coordinate (B) at (\B{1},\B{2});
  \coordinate (C) at (\C{1},\C{2});
  \coordinate (D) at (\D{1},\D{2});
  \coordinate (E) at (\E{1},\E{2});
  \coordinate (F) at (\F{1},\F{2});
  \coordinate (M) at (\M{1},\M{2});
  % \coordinate (O1) at (\Oa{1},\Oa{2});
  % \coordinate (O2) at (\Ob{1},\Ob{2});
  % \coordinate (O3) at (\Oc{1},\Oc{2});
  % \coordinate (O4) at (\Od{1},\Od{2});
  \coordinate (I1) at (\Ia{1},\Ia{2});
  \coordinate (I2) at (\Ib{1},\Ib{2});
  \coordinate (I3) at (\Ic{1},\Ic{2});
  \coordinate (I4) at (\Id{1},\Id{2});
  \coordinate (J1) at (\Ja{1},\Ja{2});
  \coordinate (J2) at (\Jb{1},\Jb{2});
  \coordinate (J3) at (\Jc{1},\Jc{2});
  \coordinate (J4) at (\Jd{1},\Jd{2});
  \coordinate (K1) at (\Ka{1},\Ka{2});
  \coordinate (K2) at (\Kb{1},\Kb{2});
  \coordinate (K3) at (\Kc{1},\Kc{2});
  \coordinate (K4) at (\Kd{1},\Kd{2});
  \coordinate (L1) at (\La{1},\La{2});
  \coordinate (L2) at (\Lb{1},\Lb{2});
  \coordinate (L3) at (\Lc{1},\Lc{2});
  \coordinate (L4) at (\Ld{1},\Ld{2});

  \draw[thick,navy] (A) -- (B) -- (D) -- (F) (A) -- (C);
  % \draw[teal] (O1) circle (\ra);
  % \draw[violet] (O2) circle (\rb);
  % \draw[purple] (O3) circle (\rc);
  % \draw[magenta] (O4) circle (\rd);
  
  \foreach \p/\placement in {A/above,B/below left,C/below right,D/below right,E/above right,F/above left,M/above right,
  I1/above,I2/above,I3/above,I4/above,
  J1/above,J2/above,J3/above,J4/above,
  K1/above,K2/above,K3/above,K4/above,
  L1/above,L2/above,L3/above,L4/above}{
    \fill[red] (\p) circle (1.5pt);
    \draw (\p) node[\placement] {$\p$};
  }

  \draw[xshift=-2.5cm,yshift=-5.5cm] node [right,text width=9cm,style=information text]
  {
    \textit{四个三角形的内心和旁心，共计有16个点。这16个点里面存在8个四点共圆，这八个圆可以分成彼此正交的两组。其中一组是互不相交的四个圆，确定一个共轴圆系，因此圆心共线，且经过点M。。}
  };

\end{tikzpicture}

\end{document}